% ---
% id: reductions/cor:register-instance
% kind: cor
% title: "The register instance: register-conditioned squashing of instances"
% summary: >-
%   When a register read by a POVM selects the branch, the branching instance
%   is a reduction instance whose pull-backs are the branches' tensored with
%   the register element. Branch entries hold at constant zero, and a
%   branch-weight entry only where (H) makes that element a multiple of the
%   identity.
% prominence: standard
% status: implemented
% uses: [thm:instrument-instance, thm:mixture, def:reduction-instance, def:emitnorm, lem:master, def:generalized-instrument]
% mentions: [rem:povmreg, cor:flag-instance]
% assumes: []
% refs: []
% provenance: newly-derived
% implements: certiqs_key.reductions.instance:combine_branches
% lean:
% notes: "The hypotheses that {N_j} is a POVM on D and, at a branch with a branch-weight entry, that N_j is the declared multiple of the identity are part of (H): decided on the declared register POVM where the submodule that declares the register declares it, and otherwise carried by that submodule's ledger; this corollary assumes nothing of its own."
% ---

A register $D$, read by a POVM that need not be projective
(Remark~\ref{rem:povmreg}), selects which branch instance acts; a branch may
be the flag instance of Corollary~\ref{cor:flag-instance}. The corollary is
the branching instance at the instrument that reads the register.

\begin{corollary}\label{cor:register-instance}
Let $\Bsys=\Bsys_0D$, let $\{N_j\}_{j\in J}$ be a POVM on the register $D$
with $J$ finite and ordered, and for each $j$ let $\mathcal{J}_j$ be a
reduction instance on $\Bsys_0$ over one common joint alphabet with classes
and weights, with handed model $\{M_{v|j}\}_{v\in\Vset}$, channel
$\Lambda_j$, transported set $\hat O^{(j)}$ and entry list $\CList_j$. Let
$\Ins_j$ be the instrument of Theorem~\ref{thm:mixture},
$\Ins_j(\rho)=\tr_{D}\big[(\idon{\Bsys_0}\otimes N_j^{1/2})\rho
(\idon{\Bsys_0}\otimes N_j^{1/2})\big]$. Then the branching instance of
Theorem~\ref{thm:instrument-instance} for $\{\Ins_j\}$ and
$\{\mathcal{J}_j\}$, the \emph{register instance}, is a reduction
instance on $\Bsys$ in both forms, with the handed model
\begin{equation}\label{eq:register-model}
  M_v=\sum_{j\in J}M_{v|j}\otimes N_j
  \quad\text{(coarse-grained)},
  \qquad
  M_{(v,j)}=M_{v|j}\otimes N_j
  \quad\text{(joint-outcome)},
\end{equation}
the channel \eqref{eq:regchannel} pinched,
$\Lambda=\PinS\circ\bigoplus_j\Lambda_j\circ\Ins_j$, the target
$\bigoplus_jF^{(j)}_v$ resp.\ $F^{(j)}_v|_j$, the hypotheses
$\bigcup_j$(H$_j$) together with ``$\{N_j\}$ is a POVM on $D$'' and, for every
$j$ in a declared set $J_{\mathrm{w}}\subseteq J$ with declared lattice
constants $q_j$, ``$N_j=q_j\idon{D}$'', and:
\begin{enumerate}[label=(\roman*),leftmargin=3.0em]
\item \emph{(transport)} for every branch $j$ and sector $s$ of
  $\mathcal{J}_j$, and for every $\hat X\in\hat O^{(j)}$,
  \begin{equation}\label{eq:register-transport}
    \Lambda^{\dagger}\big(\Pi_{(j,s)}\big)=\Lambda_j^{\dagger}(\Pi_s)\otimes N_j,
    \qquad
    \Lambda^{\dagger}\big(\Pi_j\big)=\idon{\Bsys_0}\otimes N_j,
    \qquad
    \Lambda^{\dagger}\big(\hat X|_j\big)=\Lambda_j^{\dagger}(\hat X)\otimes N_j ;
  \end{equation}
\item \emph{(the entry clause)} every entry $(\hat O,\bowtie,c)$ of a branch
  $j$ is carried as $\big(\hat O|_j-c\,\Pi_j,\bowtie,0\big)$, with pull-back
  \begin{equation}\label{eq:register-entry}
    \Lambda^{\dagger}\big(\hat O|_j-c\,\Pi_j\big)
    =\big(\Lambda_j^{\dagger}(\hat O)-c\,\idon{\Bsys_0}\big)\otimes N_j ,
  \end{equation}
  a tensor product of the branch's certified operator with $N_j\succeq0$, so
  that the certificate holds at $0$ with the branch's relation; the
  branch-weight entry $(\Pi_j,=,q_j)$ is listed for every $j\in J_{\mathrm{w}}$;
\item \emph{(the class-marginal identity)} it is discharged when every
  branch discharges it, and the flag is set in exactly that case;
\item \emph{(the declared model)} the normalizer acts branch by branch, and
  the declared model is
  \begin{equation}\label{eq:register-declared}
    \Mdec_v=\sum_{j\in J}{\Mdec}^{(j)}_v\otimes N_j
    \quad\text{resp.}\quad
    \Mdec_{(v,j)}={\Mdec}^{(j)}_v\otimes N_j ,
  \end{equation}
  with ${\Mdec}^{(j)}$ the declared model of branch $j$.
\end{enumerate}
\end{corollary}

\begin{proof}
The proof of Theorem~\ref{thm:mixture} shows that $\{\Ins_j\}$ is a
\tref{generalized instrument} (Definition~\ref{def:generalized-instrument})
with $\Ins_j^{\dagger}(X)=X\otimes N_j$ for $X\in\Lin{\Bsys_0}$, and $\Ins_j$
is given by $N_j$ alone, independently of any datum; the hypothesis that it
is an instrument is the hypothesis that $\{N_j\}$ is a POVM on $D$, the
branch-weight hypothesis of the theorem at $j\in J_{\mathrm{w}}$ is
$N_j=q_j\idon{D}$ (below), and that
$\Hs{\Bsys}$ is separable is inherited from $\Bsys_0$ and $D$.
Theorem~\ref{thm:instrument-instance} therefore applies. Its model maps are
$M_{(v,j)}[A']=\Ins_j^{\dagger}(M^{(j)}_v[A'_j])=M^{(j)}_v[A'_j]\otimes N_j$
and their sum over $j$, which at the handed data are the two models
\eqref{eq:register-model}; its channel \eqref{eq:instrument-channel} is
\eqref{eq:regchannel} composed with the sector pinching, as in
Theorem~\ref{thm:mixture}\,(i) and (iv). (i) is
\eqref{eq:instrument-transport} and the transport clause of the theorem with
$\Ins_j^{\dagger}(X)=X\otimes N_j$ substituted, the second identity being
Theorem~\ref{thm:mixture}\,(iii). (ii) is \eqref{eq:instrument-entry} with
the same substitution:
$\Ins_j^{\dagger}\big(\Lambda_j^{\dagger}(\hat O)-c\idon{\Bsys_0}\big)
=\big(\Lambda_j^{\dagger}(\hat O)-c\idon{\Bsys_0}\big)\otimes N_j$; if
$Y\succeq0$ on $\Bsys_0$ and $N_j\succeq0$ then $Y\otimes N_j\succeq0$, if
$Y\preceq0$ then $Y\otimes N_j\preceq0$, and $Y=0$ gives $0$, so the branch's
relation survives at the constant $0$. The branch weight
$\Ins_j^{\dagger}(\idon{\Bsys_0})=\idon{\Bsys_0}\otimes N_j$ is a multiple
$q_j\idon{\Bsys}$ exactly when $N_j=q_j\idon{D}$, so the theorem's
branch-weight hypothesis at $j\in J_{\mathrm{w}}$ is the hypothesis
$N_j=q_j\idon{D}$ stated here, and the entry is the theorem's (and that is
the reason Theorem~\ref{thm:mixture}\,(iii) calls the general case an
identity rather than a constraint). (iii) and (iv)
are the theorem's two closing sentences at this instrument, the declared
model \eqref{eq:register-declared} being the model map at
$A'=(A'_j)_j$ by \eqref{eq:declmodel}.
\end{proof}
